% SEGMENT OLP-0627-S01
% Part: history
% Chapter: biographies
% Section: rozsa-peter

\documentclass[../../../include/open-logic-section]{subfiles}

\begin{document}

\olfileid{his}{bio}{pet}

\olsection{ரோசா பீட்டர்}

\olphoto{peter-rozsa}{ரோசா பீட்டர்}

ரோசா பீட்டர் 1905 பிப்ரவரி 17 அன்று ஹங்கேரியின் புடாபெஸ்டில் ரோசா
பொலிட்சர் என்ற பெயரில் பிறந்தார். மீள்வரையறைச் சார்புகள் பற்றிய
பணிக்காகவே அவர் மிகவும் அறியப்படுகிறார்; மீள்வரையறைக் கோட்பாடு என்னும்
துறை உருவாவதற்கு அது இன்றியமையாததாக இருந்தது.

% SEGMENT OLP-0627-S02
கடுமையான அரசியல் சூழலில் பீட்டர் வளர்ந்தார்; பதின்ம வயதில் இருக்கும்போது
முதலாம் உலகப் போர் நடந்துகொண்டிருந்தது. இருந்தபோதிலும் புடாபெஸ்டின்
வளமான மரியா தெரேசியா பெண்கள் பள்ளியில் படிக்க முடிந்து, 1922-இல்
படிப்பை முடித்தார். பின்னர் புடாபெஸ்டின் பாஸ்மான் பீட்டர்
பல்கலைக்கழகத்தில் படித்தார்; அதற்குப் பிறகு அந்த நிறுவனம் லோராண்ட்
எட்வோஷ் பல்கலைக்கழகம் எனப் பெயர் மாற்றப்பட்டது. தந்தையின் வற்புறுத்தலால்
வேதியியல் படிக்கத் தொடங்கி, பின்னர் கணிதத்திற்கு மாறி, 1927-இல்
பட்டம் பெற்றார். மேல்நிலைப் பள்ளியில் கணிதம் கற்பிக்கத் தகுதி இருந்தாலும்,
பெருமந்தநிலை உலகப் பொருளாதாரத்தைப் பாதித்ததால் அக்காலப் பொருளாதாரச்
சூழல் கடினமாக இருந்தது. இதனால் கணிதத் தனிப்பயிற்சியாளராகவும் தனியார்
ஆசிரியராகவும் பல தற்காலிகப் பணிகளைச் செய்தார். பின்னர் கணிதத்தில்
மேற்படிப்பு மேற்கொள்ளப் பல்கலைக்கழகத்துக்குத் திரும்பினார். முதலில்
எண்கோட்பாட்டில் பணியாற்றத் திட்டமிட்டார்; ஆனால் தாம் பெற்ற முடிவுகள்
முன்பே நிறுவப்பட்டிருப்பதை அறிந்ததும் கணிதத்தையே கைவிடும் நிலைக்கு
வந்தார். கோடலின் முழுமையின்மைத் தேற்றங்களில் பணியாற்ற ஊக்குவிக்கப்பட்டார்;
அவர் அறியாமலேயே கோடலின் பல முடிவுகளை வேறு முறைகளில் நிறுவினார். இதனால்
தன்னம்பிக்கை மீண்டது. டேவிட் ஹில்பர்ட்டின் அடித்தளத் திட்டத்தால்
தூண்டப்பட்டு, மீள்வரையறைக் கோட்பாடு பற்றிய முதல் கட்டுரைகளை எழுதினார்.
1935-இல் முனைவர் பட்டம் பெற்றார்; 1937-இல் \emph{Journal of
  Symbolic Logic} இதழின் ஆசிரியர்களில் ஒருவரானார்.

% SEGMENT OLP-0627-S03
பீட்டரின் தொடக்ககாலக் கட்டுரைகள், மீள்வரையறைச் சார்புக் கோட்பாட்டை
நிறுவிய பங்களிப்புகளாகப் பரவலாகக் கருதப்படுகின்றன. \cite{Peter1935a}
இல் வெவ்வேறு வகை மீள்வரையறைகளுக்கிடையிலான உறவை ஆராய்ந்தார்.
\cite{Peter1935b} இல் குறிப்பிட்ட மீள்வரையறைச் சார்பு ஆதிம
மீள்வரையறைச் சார்பு அல்ல என்று காட்டினார். இது வில்ஹெல்ம் ஆக்கர்மனின்
முந்தைய முடிவை எளிமையாக்கியது. பீட்டரின் எளிமைப்படுத்தப்பட்ட சார்பே
இப்போது பெரும்பாலும் ஆக்கர்மன் சார்பு என்றும், சில சமயம் இன்னும்
பொருத்தமாக ஆக்கர்மன்--பீட்டர் சார்பு என்றும் அழைக்கப்படுகிறது.
மீள்வரையறைச் சார்புக் கோட்பாட்டின் முதல் நூலை அவர் எழுதினார்
\citep{Peter1951}.

% SEGMENT OLP-0627-S04
அவருடைய பணி முக்கியமானதும் தாக்கம் கொண்டதுமாக இருந்தபோதிலும்,
1945 வரையில் முழுநேரக் கற்பித்தல் பணி கிடைக்கவில்லை. இரண்டாம் உலகப்
போரின்போது ஹங்கேரியை நாசிகள் ஆக்கிரமித்திருந்த காலத்தில், யூதர்களுக்கு
எதிரான சட்டங்களால் பீட்டருக்குக் கற்பிக்க அனுமதி இல்லை. 1944-இல்
அரசாங்கம் புடாபெஸ்டில் யூதர்களுக்கான தனிமைப்படுத்தப்பட்ட குடியிருப்புப்
பகுதியை உருவாக்கியது; நகரின் மற்ற பகுதிகளிலிருந்து அது துண்டிக்கப்பட்டு,
ஆயுதமேந்திய காவலர்களால் கண்காணிக்கப்பட்டது. 1945-இல் அது விடுவிக்கப்படும்
வரை பீட்டர் அங்கே வாழக் கட்டாயப்படுத்தப்பட்டார். பின்னர் புடாபெஸ்ட்
ஆசிரியர் பயிற்சிக் கல்லூரியில் கற்பித்தார்; 1955 முதல் எட்வோஷ்
லோராண்ட் பல்கலைக்கழகத்தில் பணியாற்றினார்.
கணிதத்திற்கான Academic Doctor தகுதி பெற்ற முதல் ஹங்கேரியப் பெண்
கணிதவியலாளரும், ஹங்கேரிய
அறிவியல் கழகத்துக்குத் தேர்ந்தெடுக்கப்பட்ட முதல் பெண்ணும் அவரே.

% SEGMENT OLP-0627-S05
கணிதத்தை ஆர்வத்துடன் கற்பித்த ஆசிரியராகப் பீட்டர் அறியப்பட்டார்.
வெறுமனே விரிவுரை செய்வதைவிட மாணவர்களுடன் சேர்ந்து கணிதச் சிக்கல்களின்
இயல்பையும் அழகையும் ஆராய விரும்பினார். அதனால் மாணவர்கள் அவரை
அன்புடன் ``ரோசா அத்தை'' என்று அழைத்தனர். பீட்டர் 1977-இல்
71 வயதில் இறந்தார்.

% SEGMENT OLP-0627-S06
\begin{reading}
மேலும் வாழ்க்கை வரலாற்றுத் தகவல்களுக்கு \citep{Oconnor2014},
\citep{Andrasfai1986} ஆகியவற்றைப் பார்க்கவும். \citet{Tamassy1994}
பீட்டருடன் ஒரு சுருக்கமான நேர்காணலை நடத்தினார். கணிதம் பற்றிய
சுவாரசியமான வாசிப்புக்கு, பீட்டரின் \emph{Playing With Infinity}
நூலைப் பார்க்கவும் \citep{Peter2010}.
\end{reading}

\end{document}
